\chapter{Modulo FDCDomain}
\label{ch:fdcdomain}

\section{Responsabilità}

\module{FDCDomain} contiene i \textbf{tipi di dominio ferroviario}:
entità immutabili (o quasi), \texttt{Codable} per persistenza,
senza dipendenze da SwiftUI o da \texttt{NetworkModel}.

Il package SPM in \repo{Packages/FDCDomain/} espone un sottoinsieme
\textbf{Foundation-only} compilabile standalone; il target Xcode include
l'intera cartella \repo{FdC Railway Manager/Domain/Model/}.

\section{Grafo del dominio}

\begin{figure}[htbp]
\centering
\begin{tikzpicture}[
  type/.style={draw, ellipse, minimum width=2.4cm, minimum height=0.8cm,
    align=center, font=\footnotesize},
  arr/.style={-{Stealth[length=1.5mm]}}
]
  \node[type] (topo) {RailwayTopology};
  \node[type, below left=1cm and 0.5cm of topo] (node) {Node};
  \node[type, below right=1cm and 0.5cm of topo] (edge) {Edge};
  \node[type, below=1.8cm of node] (train) {Train};
  \node[type, below=0.6cm of train] (stop) {RelationStop};
  \node[type, right=2.2cm of edge] (seg) {TrackSegment};
  \node[type, above=0.5cm of seg] (sig) {Signal};
  \node[type, below=0.6cm of edge] (tcp) {TrackControlPoint};
  \node[type, left=1.8cm of node] (elec) {ElectrificationType};
  \node[type, right=1.8cm of node] (rc) {RoutingConstraint};

  \draw[arr] (topo) -- (node);
  \draw[arr] (topo) -- (edge);
  \draw[arr] (edge) -- (seg);
  \draw[arr] (edge) -- (tcp);
  \draw[arr] (seg) -- (sig);
  \draw[arr] (node) -- (elec);
  \draw[arr] (node) -- (rc);
  \draw[arr] (train) -- (stop);
\end{tikzpicture}
\caption{Relazioni principali tra tipi di dominio.}
\end{figure}

\section{Tipi per categoria}

\subsection{Topologia}

\begin{description}
  \item[\texttt{RailwayTopology}] Snapshot immutabile di \texttt{nodes}
    e \texttt{edges}. Delega pathfinding a \texttt{NetworkPathfinder}
    (Fase 5). Nessun \texttt{init(network:)}.
  \item[\texttt{Node}] Stazione, deposito, interscambio; vincoli di
    instradamento, altitudine, parametri Takt (\texttt{taktMinutes}).
  \item[\texttt{Edge}] Tratta tra due nodi; segmenti \texttt{TrackSegment},
    punti di controllo, elettrificazione, limite velocità.
\end{description}

\subsection{Infrastruttura di dettaglio}

\begin{description}
  \item[\texttt{TrackSegment}] Binario fisico con ordine, lunghezza,
    segnale opzionale, occupazione.
  \item[\texttt{TrackControlPoint}] Punto sulla tratta (lat/lon/alt).
  \item[\texttt{Signal}] Segnale con aspetto (\texttt{stop}, \texttt{proceed},
    \texttt{caution}).
  \item[\texttt{Switch}] Deviatione con \texttt{nodeId} e collegamenti.
  \item[\texttt{ElectrificationType}] Trazione: diesel, 3\,kV, 25\,kV, 950\,V.
  \item[\texttt{RoutingConstraint}] Vincoli di percorso per stazione.
\end{description}

\subsection{Servizio e rotazione}

\begin{description}
  \item[\texttt{Train}] Treno con fermate \texttt{RelationStop}, orari,
    assegnazione binari, \texttt{routeId}.
  \item[\texttt{RelationStop}] Fermata o transito su una stazione.
  \item[\texttt{Vehicle}] Rotazione materiale rotabile.
  \item[\texttt{TrainRoute}] Template di percorso (origine, destinazione,
    sequenza stazioni). \texttt{displayColor} in \repo{Models/TrainRoute+UI.swift}.
\end{description}

\section{Package SPM}

\begin{lstlisting}
Packages/FDCDomain/
  Package.swift
  Sources/FDCDomain/    # symlink ai .swift Foundation-only
  Tests/FDCDomainTests/
\end{lstlisting}

Tipi inclusi nel package (v0.1):
\texttt{ElectrificationType}, \texttt{Signal}, \texttt{Switch},
\texttt{RoutingConstraint}, \texttt{RelationStop}, \texttt{Vehicle},
\texttt{TrackSegment}, \texttt{TrackControlPoint}, \texttt{RailwayConstants}.

Tipi \textbf{esclusi} fino a decoupling UI / NetworkModel:
\texttt{Node}, \texttt{Edge}, \texttt{Train}, \texttt{RailwayTopology},
\texttt{LineProfilePoint}, \texttt{TrainRoute}, \texttt{Ferrovia}.

\section{Regole di evoluzione}

\begin{enumerate}
  \item Nuovi modelli → \repo{Domain/Model/}, solo \texttt{import Foundation}.
  \item Estensioni SwiftUI → file \texttt{+UI.swift} in \repo{Models/}.
  \item Rimuovere voci duplicate da \texttt{project.pbxproj} (sync group).
  \item Aggiungere test in \texttt{FDCDomainTests} quando il tipo è standalone.
\end{enumerate}
